%------------------------------------------------------------------------------%
\section{Derivadas}
\label{sec:Derivadas}

Com o limite podemos definir a derivada de uma função real que nos fornece o
valor da inclinação da reta tangente ao gráfico da função em cada ponto.

%------------------------------------------------------------------------------%
\begin{definicao}{Derivada de Função Real}{def:derivada-funcao-real}
  A \emph{Derivada}\index{Derivada} de uma função real $f(x)$ em relação
  a variável $x$ é a função $f'$ cujo valor em $x$ é
  \[
    f'(x) = \lim_{h\to 0} \frac{\,f(x+h)-f(x)\,}{h}
  \]
  desde que o limite exista.
\end{definicao}
%------------------------------------------------------------------------------%

Ao calcularmos derivadas usamos as regras de derivação listadas na proposição
a seguir.

%------------------------------------------------------------------------------%
\begin{proposicao}{Regras de Derivação}{pro:regras-derivacao}
  Sejam $f$ e $g$ funções deriváveis, $c$ e $r$ constantes, então temos que
  \begin{multicols}{2}
    \begin{enumerate}[parsep=11pt]
      \item \( \frac{d}{dx} c = 0 \)
      \item \( \frac{d}{dx} x^n = nx^{n-1} \)
      \item \( \frac{d}{dx} e^x = e^x \)
      \item \( \big(cf\big)' = cf' \)
      \item \( \big(f+g\big)' = f' + g' \)
      \item \( \big(f-g\big)' = f' - g' \)
      \item \( \big(fg\big)' = f'g + fg' \)
      \item \( \left(\frac{f}{\,g\,}\right)' = \frac{f'g - fg'}{g^2} \)
      \item Se \(h(x) = f\big(g(x)\big) \) então  \[ h'(x) = f'\big(g(x)\big)\,g'(x) \]
    \end{enumerate}
  \end{multicols}
\end{proposicao}
%------------------------------------------------------------------------------%

Uma propriedade importante de funções deriváveis é o Teorema de Rolle ou
Teorema do Valor Médio.

%------------------------------------------------------------------------------%
\begin{teorema}{Teorema de Rolle}{teo:rolle}
  Suponha\index{Teorema!de Rolle}\index{Teorema!valor médio}
  que $f$ é uma função
  \begin{enumerate}
    \item contínua  no intervalo fechado $[a, b]$
    \item derivável no intervalo aberto  $(a, b)$ e
    \item \(f(a)=f(b)\)
  \end{enumerate}
  então existe $c\in(a, b)$ onde $f'(c)=0$.
\end{teorema}
%------------------------------------------------------------------------------%

